% =====================================================================
%  7. CONCLUSION
% =====================================================================
\section{Conclusion et perspectives}\label{sec:conclusion}

Nous avons présenté TRIADE, un protocole de gouvernance du développement
logiciel pensé pour des exécutants agentiques~: trois artefacts versionnés
(intention, binaire, carte), dix règles à sanction mécanique, un verrou
croisé cryptographique, une fonction objectif lexicographique, et une
boucle de mise à jour dont le commit est gardé par la vérification.
L’article en a donné une présentation formalisée et l’a confrontée à sa
première validation à l’échelle~: le recodage complet de l’application
RoomEq — 44 spécifications avant le code, 42 modules cartographiés, 437
tests contre 228, zéro avertissement, parité fonctionnelle vérifiée,
séquence de publication incluant deux échecs d’intégration continue honorés
— avec une technique au passage~: les spécifications rétro-actives de
l’installateur scellé. Les menaces sur la validité ont été traitées comme
une contribution à part entière~: cas unique sans témoin,
auto-vérification par l’agent auteur, et écart assumé entre la
refactorisation source démontrée et la génération binaire native promise.

Trois perspectives découlent directement de ces limites.
Premièrement, la \emph{réplication}~: appliquer le protocole à d’autres
projets, piles et agents, idéalement par des équipes indépendantes, pour
convertir la plausibilité en généralité. Deuxièmement, la
\emph{réalisation}~: franchir les jalons M0–M7 de la feuille de route —
chaîne notation-octets-exécutable sans compilateur, patch à chaud sur code
en exécution, portabilité inter-OS puis inter-ISA — qui seuls valideront
le cœur natif du protocole. Troisièmement, l’\emph{instrumentation}~:
mesurer ce que l’étude a seulement ciblé — coût réel en tokens par
édition, couverture de code, coût total de la discipline comparé à un
développement non gouverné. Le fil conducteur reste celui de
l’introduction~: développer avec un agent exige de l’agent qu’il prouve
ce qu’il affirme, et du projet qu’il survive à l’agent. La démonstration
RoomEq établit que la première moitié de cette exigence est tenable à
l’échelle d’une application réelle~; la seconde est, par construction, en
travail permanent.
